% answers_short.tex —— 1995 年·上海卷·理工农医类 · 简短答案速查
% 由 tools/make_short_answers.py 自 answers.tex 的 \gksolhead/\gkanswer 提取，
% 并与 1995/1995_数学真题_上海理_答案要点.md 逐题比对后生成，勿手改（改源 MD 后重跑）。

\begin{gknotice}
\textbf{使用说明}：本节只给最终答案便于快速核对，完整推导见解析版。答案为固定来源 PDF 第 5—9 页卷内「答案及评分标准」所载（出版者/版次未核实），非官方标准答案；第 24 题按题设补正 $m=n$ 情形并保留原答案。\textbf{非官方标准答案}。
\end{gknotice}

\gkpart{选择题}
\begin{center}\begin{tabular}{ccccccccc}
\toprule
题号 & 1 & 2 & 3 & 4 & 5 & 6 & 7 & 8 \\
\midrule
答案 & C & B & B & C & A & A & D & B \\
\bottomrule
\end{tabular}\end{center}

\gkpart{填空题}
\begin{itemize}[leftmargin=2.6em]
  \item[\textbf{9.}] $x<-3$ 或 $x>4$。
  \item[\textbf{10.}] $3x+4y=0$ 和 $3x-4y=0$。
  \item[\textbf{11.}] $y=54.8\left(1+\frac{x}{100}\right)^8$。
  \item[\textbf{12.}] $1$。
  \item[\textbf{13.}] $11$。
  \item[\textbf{14.}] $y^2=-8x+8$。
  \item[\textbf{15.}] $x^2+(y-1)^2=1$。
  \item[\textbf{16.}] $3x^2-y^2=1$。
  \item[\textbf{17.}] $\left[-\frac{3\pi}{2},\frac\pi2\right]$（来源区间记号；原题定义域为开区间，端点处导数零，闭区间上严格递增）。
  \item[\textbf{18.}] $A_1$ 组 $-2<r<0$；$B_1$ 组 $e$。
  \item[\textbf{19.}] $A_1$ 组 $350$ 种；$B_1$ 组 $\frac{25}{33}$（或 $0.76$）。
  \item[\textbf{20.}] $A_1$ 组 $37.70\,\mathrm{dm}^3$（原评分接受 $37.68$ 或 $37.69$）；$B_1$ 组 $2+i$。
\end{itemize}

\gkpart{解答题}
\begin{itemize}[leftmargin=2.6em]
  \item[\textbf{21.}] $A_2$ 组 $-\frac45$；$B_2$ 组（1）$\overrightarrow{OD}=\left(0,-\frac12,\frac{\sqrt3}{2}\right)$；（2）$\cos\theta=-\frac{\sqrt{10}}5$。
  \item[\textbf{22.}] $A_2$ 组 $(z_1,z_2)=\left(-\frac12+\frac{\sqrt3}{2}i,1\right)$ 或二者互换；$B_2$ 组（1）$f(x)=(x+1)^2$；（2）$t=1-2^{-1/3}$。
  \item[\textbf{23.}] （1）$V=2\sqrt3$；（2）二面角大小为 $\arctan\frac{2\sqrt3}{3}$。
  \item[\textbf{24.}] （1）$S=\dfrac{4m^2n^2\tan\theta}{n^2+m^2\tan^2\theta}$；（2）$n\le m$ 时 $u=2mn$（含 $m=n$，按题设补正），$m<n$ 时 $u=\dfrac{4m^2n^2}{m^2+n^2}$；（3）$\dfrac mn\in(2-\sqrt3,+\infty)$（卷内原答案作 $(2-\sqrt3,1)\cup(1,+\infty)$，漏 $m=n$；按题设补全并保留原文）。
  \item[\textbf{25.}] （1）$f(x)=x^2-(t+2)x+t+1$；（2）$a_n=\dfrac{(t+1)^{n+1}-1}{t}$，$b_n=\dfrac{t+1}{t}\left[1-(t+1)^n\right]$；（3）$r_n=\dfrac{\sqrt2(t+1)^{n+1}}{t+2}$，$S_n=\dfrac{2\pi(t+1)^4\left((t+1)^{2n}-1\right)}{t(t+2)^3}$。
\end{itemize}
